% Part: computability
% Chapter: recursive-functions
% Section: primitive-recursion

\documentclass[../../../include/open-logic-section]{subfiles}

\begin{document}

\olfileid{cmp}{rec}{pre}
\olsection{आदिम पुनरावर्तन}

नैसर्गिक संख्यांचे एक वैशिष्ट्य असे आहे की, उत्तरवर्ती क्रिया~$+1$
सांत वेळा लागू करून प्रत्येक नैसर्गिक संख्या~$0$ पासून गाठता येते---कोणतीही
नैसर्गिक संख्या एकतर~$0$ असते किंवा~$0$ च्या उत्तरवर्तीचा \dots{}
उत्तरवर्ती असते. या तथ्याचा उपयोग करणारे फलन~$h\colon\Nat \to \Nat$
विनिर्दिष्ट करण्याचा एक मार्ग असा: (a)~आदान~$0$ साठी $h$ चे मूल्य काय
आहे ते विनिर्दिष्ट करणे, आणि (b)~$h(x)$ चे मूल्य दिले असता $h(x+1)$ चे
मूल्य कसे संगणित करायचे तेही विनिर्दिष्ट करणे. कारण (a) आपल्याला
$h(0)$ काय आहे ते थेट सांगते, त्यामुळे~$0$ साठी $h$ परिभाषित होते.
आता (b) ने दिलेली सूचना $x=0$ साठी वापरून आपण~$h(0)$ वरून
$h(1) = h(0+1)$ संगणित करू शकतो. तीच सूचना $x=1$ साठी वापरून
आपण~$h(1)$ वरून $h(2) = h(1+1)$ संगणित करतो, आणि असेच पुढे.
प्रत्येक नैसर्गिक संख्या~$x$ साठी, आपण अखेरीस $h(x)$ वरून $h(x+1)$
परिभाषित करतो त्या पायरीपर्यंत पोहोचतो; आधारमूल्यापासून अशी पावले
घेतल्यामुळे सर्व $x \in \Nat$ साठी $h(x)$ परिभाषित होते.

उदाहरणार्थ, पुढील दोन समीकरणांनी आपण $h\colon \Nat \to \Nat$ विनिर्दिष्ट
करतो असे समजा:
\begin{align*}
h(0) & =  1\\
h(x+1) & =  2 \cdot h(x)
\end{align*}
गुणाकार कसा करायचा हे आपल्याला आधीच माहीत असेल, तर वरील (a) आणि~(b)
साठी आवश्यक माहिती ही समीकरणे देतात. दुसरे समीकरण क्रमाक्रमाने लागू
केल्यावर आपल्याला पुढील गोष्टी मिळतात:
\begin{align*}
  h(1) & = 2\cdot h(0) = 2,\\
  h(2) & = 2\cdot h(1) = 2\cdot 2,\\
  h(3) & = 2 \cdot h(2) = 2\cdot 2 \cdot 2,\\
  & \vdots
\end{align*}
आपण विनिर्दिष्ट केलेले फलन~$h$ हे $h(x) = 2^x$ आहे, हे दिसते.

नैसर्गिक संख्यांच्या वैशिष्ट्यपूर्ण गुणधर्मामुळे हे दोन निकष पूर्ण करणारे
एकमेव फलन~$h$ असते. अशा समीकरणांच्या जोडीला फलन~$h$ ची
\emph{आदिम पुनरावर्तनाने केलेली व्याख्या} म्हणतात. तिला तसे म्हणण्याचे
कारण म्हणजे आपण~$h$ ची व्याख्या ``पुनरावर्ती रीतीने'' करतो, म्हणजे
व्याख्येत, विशेषतः दुसऱ्या समीकरणाच्या उजव्या बाजूला, $h$ स्वतःच येते.
ती ``आदिम'' असते, कारण $h(x+1)$ परिभाषित करताना आपण केवळ~$h(x)$ चे
मूल्य, म्हणजे अगदी आधीचे मूल्य, वापरतो. फलनाची~$\Nat$ वर पुनरावर्ती
रीतीने व्याख्या करण्याचा हा सर्वांत सोपा मार्ग आहे.

बेरीज आणि गुणाकार यांसारखी आणखी मूलभूत फलनेही आपण आदिम पुनरावर्तनाने
परिभाषित करू शकतो. मात्र या प्रकरणांत संबंधित फलने $2$-स्थानी असतात.
आपण आदानस्थानांपैकी एक स्थिर ठेवतो आणि दुसरे पुनरावर्तनासाठी वापरतो.
उदाहरणार्थ, $\Add(x, y)$ परिभाषित करण्यासाठी आपण~$x$ स्थिर ठेवून आधी
$y=0$ साठी मूल्य परिभाषित करू शकतो आणि नंतर~$y$ च्या साहाय्याने $y+1$
साठीचे मूल्य परिभाषित करू शकतो. $x$ स्थिर असल्यामुळे ते परिभाषक
समीकरणांच्या डाव्या आणि उजव्या दोन्ही बाजूंना येईल.
\begin{align*}
\Add(x,0) & =  x\\
\Add(x,y+1) & =  \Add(x,y)+1
\end{align*}
ही समीकरणे सर्व $x$ \emph{आणि}~$y$ साठी $\Add$ चे मूल्य विनिर्दिष्ट
करतात. उदाहरणार्थ, $\Add(2,3)$ शोधण्यासाठी आपण $x = 2$ साठी परिभाषक
समीकरणे लागू करतो: पहिल्या समीकरणाने $\Add(2,0) = 2$ शोधतो, आणि नंतर
दुसरे समीकरण क्रमाक्रमाने वापरून $\Add(2,1) = 2 + 1 = 3$,
$\Add(2, 2) = 3 + 1 = 4$, $\Add(2, 3) = 4 + 1 = 5$ शोधतो.

$\Add$ च्या व्याख्येत आपण दुसऱ्या समीकरणाच्या उजव्या बाजूला $+$ वापरले,
पण केवळ~$1$ ची भर घालण्यासाठी. दुसऱ्या शब्दांत, आपण उत्तरवर्ती फलन
$\Succ(z) = z+1$ वापरले आणि $\Add(x,y+1)$ परिभाषित करण्यासाठी ते मागील
मूल्य~$\Add(x,y)$ ला लागू केले. म्हणून मागील मूल्याला लागू केलेल्या एका
फलनाच्या रूपात पुनरावर्ती व्याख्येचा विचार करता येतो. परंतु त्या फलनाला
केवळ मागील मूल्यावरच नव्हे, तर $x$ आणि~$y$ यांवरही अवलंबून राहण्याची
मुभा दिल्याने काही बिघडत नाही---आणि कधीकधी ते आवश्यक असते. पुढील
समीकरणे पाहा:
\begin{align*}
  \Mult(x,0) & =  0 \\
  \Mult(x,y+1) & =  \Add(\Mult(x,y),x)
\end{align*}
ही, आधीचे मूल्य $\Mult(x,y)$ आणि पहिले आदान~$x$ या दोहोंना फलन
$\Add$ लागू करून मिळणारी, फलन $\Mult$ ची आदिम पुनरावर्ती व्याख्या आहे.
ती सर्व आदाने $x$ आणि~$y$ यांच्यासाठी फलन~$\Mult(x,y)$ परिभाषितही करते.
उदाहरणार्थ, $\Mult(2,3)$ हे $\Mult(2,0)$, $\Mult(2,1)$,
$\Mult(2,2)$ आणि~$\Mult(2,3)$ क्रमाक्रमाने संगणित करून निर्धारित होते:
\begin{align*}
  \Mult(2,0) & = 0\\
  \Mult(2,1) & = \Mult(2,0+1) =
  \Add(\Mult(2,0), 2) = \Add(0, 2) = 2\\
  \Mult(2,2) & = \Mult(2,1+1) =
  \Add(\Mult(2,1), 2) = \Add(2, 2) = 4\\
  \Mult(2,3) & = \Mult(2,2+1) =
  \Add(\Mult(2,2), 2) = \Add(4, 2) = 6
\end{align*}

मग सर्वसाधारण आकृतिबंध असा आहे: फलन~$h(x_0, \dots, x_{k-1}, y)$ ची
आदिम पुनरावर्ती व्याख्या देण्यासाठी आपण दोन समीकरणे देतो. पहिले समीकरण
फलन~$h$ चा उल्लेख न करता $h(x_0, \dots, x_{k-1}, 0)$ चे मूल्य
परिभाषित करते. दुसरे समीकरण $h(x_0, \dots, x_{k-1}, y)$, इतर आदाने
$x_0$, \dots,~$x_{k-1}$ आणि~$y$ यांच्या साहाय्याने $h(x_0,
\dots, x_{k-1}, y+1)$ चे मूल्य परिभाषित करते. त्या दुसऱ्या समीकरणात
केवळ~$h$ चे अगदी आधीचे मूल्य वापरता येते. या दोन समीकरणांच्या उजव्या
बाजूंनी दिलेल्या क्रियाच स्वतः फलने $f$ आणि~$g$ आहेत असे मानल्यास,
आदिम पुनरावर्तनाने नवीन फलन~$h$ परिभाषित करण्याचा सर्वसाधारण आकृतिबंध
असा आहे:
\begin{align*}
  h(x_0, \dots, x_{k-1}, 0) & = f(x_0, \dots, x_{k-1})\\
  h(x_0, \dots, x_{k-1}, y+1) & =
  g(x_0, \dots, x_{k-1}, y, h(x_0, \dots, x_{k-1}, y))
\end{align*}
$\Add$ च्या बाबतीत $k=1$ आणि $f(x_0) = x_0$ (एकरूपता फलन) असते,
तसेच $g(x_0, y, z) = z + 1$ (आपल्या तिसऱ्या आदानाचा उत्तरवर्ती
परत करणारे $3$-स्थानी फलन) असते:
\begin{align*}
  \Add(x_0, 0) & = f(x_0) = x_0\\
  \Add(x_0, y+1) & = g(x_0, y, \Add(x_0, y)) =
  \Succ(\Add(x_0, y))
\end{align*}
$\Mult$ च्या बाबतीत $f(x_0) = 0$ (नेहमी~$0$ परत करणारे स्थिर फलन)
आणि $g(x_0, y, z) = \Add(z,x_0)$ (आपल्या शेवटच्या व पहिल्या आदानांची
बेरीज परत करणारे $3$-स्थानी फलन) असते:
\begin{align*}
  \Mult(x_0,0) & =  f(x_0) = 0 \\
  \Mult(x_0,y+1) & =  g(x_0, y, \Mult(x_0,y)) =
  \Add(\Mult(x_0,y), x_0)
\end{align*}

\end{document}
